\documentclass[../main.tex]{subfiles}
\begin{document}
\chapter{Applications of Linear Programs}
\section{Game Theory}
\subsection{Two Person Games}
Consider a game with two players where $P_1$ can pick from any of $\{1, \ldots, m\}$ actions and $P_2$ can pick from any of $\{1, \ldots, n\}$ actions.
\begin{definition}[Payoff Matrix]
  The \textit{payoff matrix} $A$ for $P_1$ is an $m \times n$ matrix with entries $A_{i j}$, where $A_{i j}$ is the score attained by $P_1$ if $P_1$ plays action $i$ and $P_2$ plays action $j$.
\end{definition}
\begin{example}
  For a game of rock paper scissors, each player can pick either rock (call this 1), paper (call this 2), or scissors (call this 3).
  The payoff matrix for $P_1$ is:
  \[
    A = \begin{pmatrix}
    0 & -1 & 1 \\
    1 & 0 & -1 \\
    -1 & 1 & 0 \\
    \end{pmatrix}
  \]
\end{example}
If $P_1$ plays $i$ then $P_1$ can expect to have payoff:
\[
  i \mapsto \min_{j \in \{1, \ldots, n\}} A_{i j}
\]
$P_1$ now tries to solve:
\[
  \max_{i \in \{1, \ldots, m\}} \min_{j \in \{1, \ldots, n\}} A_{i j}
\]
and $P_2$ will try to solve:
\[
  \min_{j \in \{1, \ldots, n\}} \max_{i \in \{1, \ldots, m\}} A_{i j}
\]
\begin{example}
  Consider the payoff matrix:
  \[
    A = \begin{pmatrix}
    1 & 2 \\
    3 & 4 \\
    \end{pmatrix}
  \]
  For this payoff matrix, $P_1$ decides they should play $i = 2$ and $P_2$ decides they should play $j = 1$ and so both their decisions coincide on the cell $A_{2 1} = 3$.

  Now consider the payoff matrix:
  \[
    A' = \begin{pmatrix}
    4 & 2 \\
    1 & 3 \\
    \end{pmatrix}
  \]
  In this case $P_1$ will play $i = 2$ and then will expect $P_2$ to choose row $j = 1$.
  However $P_2$ will chose $j = 2$ and expect $P_1$ to choose $i = 2$.
  Unlike with $A$, the choice of the two players does not coincide.
  In games where players alternate turns, this is ok, however in games where both players play simultaneously, the players end up stuck in a loop.
\end{example}
\begin{definition}[Strategy]
  A \textit{strategy} is a probability distribution for the actions that a player picks.
\end{definition}
\begin{definition}[Pure and Mixed Strategies]
  A \textit{pure strategy} is a strategy where one single action has probability 1, i.e. the player will always pick a particular action.
  If a strategy is not pure, then we call it a \textit{mixed strategy}.
\end{definition}
Let the strategy for $P_1$ be $i$ with probability $p_i$ and let the strategy for $P_2$ be $j$ with probability $q_i$.

$P_1$'s expected score is this then:
\[
  \sum_{i = 1}^{m} p_i A_{i j}
\]
So for optimal play, $P_1$ wants to adopt the strategy $\vec{p} = (p_1, \ldots, p_m)^{\trans}$ that maximises:
\[
  \max_{\vec{p} \in P_m} \left(\min_{j \in \{1, \ldots, n\}} \sum_{i = 1}^{m} p_i A_{i j}\right)
\]
where $P_m$ is the set of discrete probability distributions for $\{1, \ldots, m\}$.
\subsection{Linear Programming for Strategies}
We can express $P_1$'s optimisation problem as:
\begin{maxi}
  {\R}{v}{}{}
  \addConstraint{A^{\trans} \vec{p}}{\geq v e}
  \addConstraint{\vec{e}^{\trans} \vec{p}}{= 1}
  \addConstraint{\vec{p}}{\geq \vec{0}}
\end{maxi}
where $v$ will be the maximal score for $P_1$ and $\vec{e} = (1, \ldots, 1)^{\trans} \in \R^{m}$.

Similarly for $P_2$, we can express their optimisation problem as:
\begin{mini*}
  {\R}{w}{}{}
  \addConstraint{A\vec{q}}{\leq w \vec{e}'}
  \addConstraint{\vec{e}'^{\trans} \vec{p}}{= 1}
  \addConstraint{\vec{q}}{\geq \vec{0}}
\end{mini*}
where $w$ will be the maximal score for $P_2$ and $\vec{e}' = (1, \ldots, 1)^{\trans} \in \R^{n}$

These two problems actually turn out to be dual to one another.
To find the dual of $P_2$'s problem, we find the Lagrangian:
\[
  L(w, \vec{q}, \vec{s}, \vec{\lambda}_1, \lambda_2) = w + \vec{\lambda}^{\trans}_1(A\vec{q} + \vec{s} - w \vec{e}) - \lambda_2(\vec{e}^{\trans} \vec{q} - 1)
\]
The set $\Lambda$ is then:
\[
  \Lambda = \{\vec{\lambda}_1: \vec{\lambda}^{\trans}_1 \vec{e} = 1,\ \vec{\lambda}^{\trans}_1 A \geq \lambda_2\vec{e}^{\trans}, \vec{\lambda_1} \geq \vec{0}\}
\]
and then when $\vec{\lambda} \in \Lambda$, $\min L = \lambda_2$.
Thus the dual of the problem is:
\begin{maxi*}
  {\R}{\lambda_2}{}{}
  \addConstraint{\vec{\lambda}^{\trans}_1 \vec{e}}{= 1}
  \addConstraint{\vec{\lambda}_1}{\geq \vec{0}}
  \addConstraint{A^{\trans} \vec{\lambda_1}}{\geq \lambda_2 \vec{e}'}
\end{maxi*}
Upon relabelling $\lambda_2$ by $v$ and $\vec{\lambda}_1$ by $\vec{p}$ we obtain $P_1$'s optimisation problem.
\begin{theorem}[Optimal Strategies]
  A strategy $\vec{p}$ is optimal for $P_1$ if and only if there exists $\vec{q}$ and $v$ such that:
  \begin{enumerate}
    \item $A^{\trans} \vec{p} \geq v \vec{e}$, $\vec{e}^{\trans} p = 1$, $\vec{p} \geq \vec{0}$ (Primal feasibility)
    \item $A\vec{q} \leq v \vec{e}$, $\vec{e}^{\trans} \vec{q} = 1$, $\vec{q} \geq \vec{0}$ (Dual feasibility)
    \item $v = \vec{p}^{\trans} A \vec{q}$ (Complimentary slackness)
  \end{enumerate}
\end{theorem}
\begin{proof}
  We need only check the complimentary slackness equations.
  $(\vec{p}, v)$ and $(\vec{q}, w)$ are primal/dual optimal if and only if both:
  \[
    (A \vec{q} - w \vec{e})^{\trans} \vec{p} = 0 \text{ and } \vec{q}^{\trans}(A^{\trans} \vec{p} - v\vec{e}) = 0
  \]
  which is equivalent to $v = w = \vec{p}^{\trans} A \vec{q}$.
\end{proof}
\end{document}
